\subsection{The 1620-arrow realization and its complete relation system}
\label{sec:1620-audit}

\begin{result}{B.5}{Proved}{The specified realization of EXIT\_TOWER B.25}
The exact path realization of EXIT\_TOWER B.21--B.25 has $542$ objects
and $1620$ oriented arrows.  On its specified $541$-edge spanning tree,
the $1079$ based non-tree cycles have the complete system of $1080$
defining relations recorded in \texttt{check\_fivepoint\_full.json}.
The comparison to the fundamental groupoid of $M_{0,5}$ is an equivalence.
The original diagram and arrow count are in
\cite[manuscript folios~30 and~32, PDF scans~31 and~33]{GrothendieckCote148}.
The scans confirm the lower fractional twist as $1/4$; the
first-pass transcription's $1/3$ is corrected at that occurrence.
Every arrow word, every tree transport and all $1080$ expanded relations
have been independently reconstructed and checked.  This confirms B.25
with its stated choices of paths and pins.  The precise comparison to
an additional manuscript path convention is the change-of-generators
formula below; such a convention is not supplied by an incidence graph
alone.
\end{result}

\begin{proof}
We use labels $0,1,2,3,4$, permutation composition
$(gh)(i)=g(h(i))$, and the lexicographically least representative
of each right coset.  Set
\[
 \begin{gathered}
 H=\langle(12)(34)\rangle,\qquad
 H_B=\langle(012),(12)(34)\rangle,\\
 H_Q=\langle(1324)\rangle,
 \quad H_P=\langle(03124),(12)(34)\rangle.
 \end{gathered}
\]
Their orders are $2,6,4,10$.  The sixty core incidences are
$gH\mapsto(gH_B,gH_Q)$, with twenty bipyramids and thirty square
pyramids.  The graph is connected.  The rest of the specified objects
and arrows have the following counts; these are counts of labelled
objects with the stipulated markings.
\[
\begin{array}{c|rrrrrrrrr|r}
\text{objects}&B&Q&P&V_R&V_L&T_R&T_M&T_L&W&\text{total}\\ \hline
\text{number}&20&30&12&60&60&60&120&60&120&542
\end{array}
\]
\[
\begin{array}{l|r}
\text{arrow family}&\text{number}\\ \hline
 B\longrightarrow Q\text{ core}&60\\
 V_R\longrightarrow P\text{ real erasures}&60\\
 V_R\longrightarrow V_L\text{ quarter turns}&60\\
 V_L\longrightarrow B\text{ imaginary erasures}&60\\
 V_L\longrightarrow Q\text{ specified composites}&120\\
 T_R\longrightarrow T_M, T_R\longrightarrow V_R&120+120\\
 T_M\longrightarrow T_L, T_M\longrightarrow V_L&120+120\\
 T_L\longrightarrow V_R, T_L\longrightarrow Q&120+60\\
 W\longrightarrow W, W\longrightarrow B&120+120\\
 T_M\longrightarrow W\text{ with three pin choices}&360\\ \hline
 \text{total}&1620
\end{array}                                                     \tag{B5.1}
\]

We recall the actual path data so that the presentation is reproducible.
Put $\epsilon=1/10$ and
$V(q)=(\infty,-q,q,-1,1)$.  The core is
$V(ir)$ for $1/\sqrt3\leq r\leq1$.  The right and left one-node
objects are $V(\epsilon)$ and $V(i\epsilon)$.
The quarter turn is the segment from $\epsilon$ to $i\epsilon$;
radial interpolation identifies it with the positive circular quarter
turn without meeting $0,1,-1$.  The two erasures use the intervals
$\epsilon\to\sqrt5-2$ and $i\epsilon\to i/\sqrt3$.
Each of the $120$ direct $f_{\vec Q,a}$ paths is, by the convention of
B.22, the imaginary erasure followed by one of the two other core
arrows at the resulting bipyramid.  This fixes its relative path class.
The pentagon endpoint identifications have been checked exactly over
$\mathbb Q(\sqrt5)$ for all permutations.  Its twelve normalized
configurations are distinct; all real erasures preserve their real
order and hence have no projected braid crossing.

The second tier starts at
\[
\begin{split}
 T_R&=(\infty,-1-\epsilon,1+\epsilon,-1+\epsilon,1-\epsilon),\\
 T_M&=(\infty,-1-i\epsilon,1+\epsilon,-1+i\epsilon,1-\epsilon),\\
 T_L&=(\infty,-1-i\epsilon,1+i\epsilon,-1+i\epsilon,1-i\epsilon).
\end{split}
\]
The six representative paths are the affine segments
$T_R\to T_M$, $T_R\to R_*$, $T_M\to T_L$,
$T_M\to M_*$, $T_L\to L_*$, and $T_L\to L_{**}$, where
\[
\begin{split}
 R_*&=(\infty,-(1-\epsilon)^{-1},0,-(1+\epsilon)^{-1},-1/2),\\
 M_*&=(\infty,-(1-i\epsilon)^{-1},0,-(1+i\epsilon)^{-1},-1/2),\\
 L_*&=(\infty,-i\epsilon,i,i\epsilon,-i),\\
 L_{**}&=(\infty,-1-i,1+i,-1+i,1-i).
\end{split}
\]
Their endpoint identifications are respectively $T_M$, $p_RV_R$,
$T_L$, $p_RV_L$, $p_LV_R$, $p_QQ$, with
\[
 p_R=(2,1,3,0,4),\quad p_L=(0,1,3,4,2),\quad p_Q=(0,4,3,1,2).
\]
The first five families use all $120$ permutations, retaining the
appropriate source or target cosets; the last uses sixty cosets.
The exact normalized endpoints of every relabelling have been checked.

For the final tier set $\zeta=(-1+i\sqrt3)/2$,
$\xi=(1+i\sqrt3)/2$ and
\[
 W(q)=(1,\zeta,\zeta^2,-q^{-1},q^{-1}),\qquad W=W(\epsilon).
\]
The sixth turn uses $q(t)=\epsilon(1-t+\xi t)$ and has endpoint
$gp_\tau$, where $p_\tau=(1,2,0,4,3)=(012)(34)$.
Its six translates concatenate to a positive full normal turn.
The erasure is
\[
 (1,\zeta,\zeta^2,-(1-t)/\epsilon,
                         1/(\epsilon(1-t))),
                     \qquad0\leq t\leq1,
\]
with its projective endpoint at $t=1$; that endpoint is the bipyramid.
For the three paths from $T_M$, use
\[
 C(\lambda)=(0,\lambda,1,\overline\lambda,\infty),\qquad
 \lambda_*=\frac{-19+i}{181},
\]
and the segments from $\lambda_*$ to
\[
 \lambda_0=-\frac{30}{91}+\frac{110\sqrt3}{819}i,
 \quad\lambda_1=\frac{27}{37}+\frac{30\sqrt3}{407}i,
 \quad\lambda_2=\frac{121}{40}+\frac{33\sqrt3}{40}i.
\]
Their endpoints are $gp_jW$, with $p_j=(024)^jp_0$ and
$p_0=(0,1,3,4,2)$.  All three segments remain in the upper half-plane,
so their configurations remain distinct.  These are the exact three
pin choices of B.25.  The path $j=0$ stays inside a disc separating
its moving triple from the fixed marks $1,\infty$; it is the stipulated
one-node path used to attach the new vertices to the tree.

Here is the exact braid computation.  Normalize labels $0,1,4$ to
$0,1,\infty$ and use base configuration $(0,1,10,20)$ for the four
finite coordinates.  To reach a vertex, move coordinate two and then
coordinate three through the detour point $N+iN^2$, taking the least
integer $N\geq2$ for which the two segments avoid all stationary
coordinates.  This specifies every transport without a numerical
choice.  A pentagon uses its first incoming real-erasure transport.
Project a complex coordinate $z$ to
$\operatorname{Re}z+\operatorname{Im}z/37$ and use transverse
coordinate $\operatorname{Im}z-\operatorname{Re}z/37$.
For a linear segment, every crossing time is obtained exactly by
solving a linear equation over $\mathbb Q(\sqrt3)$.
The sign is positive when the previously left strand is below the
other strand in this transverse coordinate.  The computation checks
that all events are nondegenerate and involve adjacent projected
strands, and records the resulting word in $\sigma_1,\sigma_2,\sigma_3$.

Straightening the normalized rational paths is justified by explicit
certificates, not sampling.  Write each original projective coordinate
as a pair of degree-at-most-one homogeneous polynomials and let
$D_{ij}$ be their determinant.  On a rational interval $[a,b]$, each
$D_{ij}$ has degree at most two.  Its three Bernstein coefficients are
$D(a)$, $2D((a+b)/2)-(D(a)+D(b))/2$, $D(b)$.
The certificate checks for each such coefficient $z$, with midpoint $m$,
\[
 \operatorname{Re}(z\overline{D_{ij}(m)})>0,
 \qquad
 |\operatorname{Im}(z\overline{D_{ij}(m)})|
       <(2-\sqrt3)\operatorname{Re}(z\overline{D_{ij}(m)}).
                                                               \tag{B5.2}
\]
Convexity of the Bernstein representation keeps every determinant in
that open cone.  A normalized coordinate difference is, up to constant
sign,
$D_{ij}D_{ac}D_{bc}/(D_{ic}D_{jc}D_{ba})$ for normalizing labels
$a,b,c$.  Its argument therefore remains within $6\pi/12=\pi/2$
of its midpoint ray.  The true difference and its endpoint chord are
in the same open half-plane; their linear homotopy never vanishes.
This proves collision-free straightening simultaneously for every
relabelling.  All inequalities in (B5.2) are decided by rational
arithmetic and comparison in $\mathbb Q(\sqrt3)$.
For the five final-tier families the exact subdivision counts are
$4,8,3,12,16$, agreeing with B.25.  All certificates and endpoints are
stored with the arrow words.

We then comb each pure braid using the exact forgetting fibration of
configuration spaces.  With
\[
 (a,b,x,y,z)=(A_{13},A_{23},A_{14},A_{24},A_{34}),\qquad
 A_{ij}=\sigma_{j-1}\cdots\sigma_{i+1}\sigma_i^2
                \sigma_{i+1}^{-1}\cdots\sigma_{j-1}^{-1},
\]
the resulting group is $F(x,y,z)\rtimes F(a,b)$.
The chosen normalization keeps the first two points fixed, so the
scalar linking number is zero and no $A_{12}$ remains.
For the fibre step, delete the last strand from the braid $w$ to obtain
$w_0$, and consider $u=w_0^{-1}w$.
In the Artin action write $X_m^u=C X_m C^{-1}$ by removing matching
inverse initial and terminal letters.  Deleting $X_m$ from $C$ gives
the word in the free point-pushing generators $A_{im}$.  This is an
isomorphism on that free kernel: it takes its explicit free basis to
the free meridian basis.  Recursion therefore gives the braid itself.
All $1620$ resulting words have been checked again in the faithful
Artin action on $F_4$, using
$\sigma_i:X_i\mapsto X_iX_{i+1}X_i^{-1}$,
$X_{i+1}\mapsto X_i$.
Faithfulness follows by fixing the disjoint arcs from a boundary
basepoint to the punctures and cutting to a disc; an automorphism
fixing the peripheral basis preserves those arcs up to isotopy, and
the remaining disc homeomorphism has the Alexander isotopy.

For the core, the breadth-first tree of B.21 has forty-nine edges.
The eleven non-tree edge numbers are
\[
                 2,3,6,8,9,13,19,22,42,45,48.
\]
The actual normalized path calculation gives exactly
\[
\begin{array}{c|l@{\qquad}c|l}
 g_1&abyz&g_2&z\\
 g_3&abz&g_4&z^{-1}y^{-1}z\\
 g_5&axz&g_6&z^{-1}x^{-1}z\\
 g_7&z^{-1}xz&g_8&b^{-1}z^{-1}y^{-1}\\
 g_9&b^{-1}a^{-1}z^{-1}y^{-1}x^{-1}&g_{10}&abxyz\\
 g_{11}&a&&
\end{array}                                                     \tag{B5.3}
\]
All eleven words have been verified as equal braids, and not just in
abelianization.  The inverse dictionary is
\[
 a=g_{11},\quad b=g_{11}^{-1}g_3g_2^{-1},\quad
 x=g_2g_6^{-1}g_2^{-1},\quad
 y=g_2g_4^{-1}g_2^{-1},\quad z=g_2.                          \tag{B5.4}
\]
The six elimination relators are
\[
\begin{aligned}
 g_1&=g_3g_4^{-1},&g_5&=g_{11}g_2g_6^{-1},&g_7&=g_6^{-1},\\
 g_8&=g_2g_3^{-1}g_{11}g_4g_2^{-1},&
 g_9&=g_2g_3^{-1}g_4g_6g_2^{-1},\\
 g_{10}&=g_3g_6^{-1}g_4^{-1}.
\end{aligned}                                                   \tag{B5.5}
\]
For the other six put $c(u,v)=v^{-1}u^{-1}vu$ and substitute (B5.4)
in
\[
\begin{array}{lll}
 axa^{-1}=z^{-1}xz,&aya^{-1}=c(x,z)yc(x,z)^{-1},&aza^{-1}=c(x,z)z,\\
 bxb^{-1}=x,&byb^{-1}=z^{-1}yz,&bzb^{-1}=c(y,z)z.
\end{array}                                                     \tag{B5.6}
\]
The forgetting fibration has fibre a three-punctured plane and base
$M_{0,4}$.  Its base has zero $\pi_2$ and free fundamental group;
its exact sequence therefore splits after the chosen lifts of $a,b$.
Conjugation on the fibre generators gives exactly (B5.6), also verified
by the Artin action.  Thus (B5.6) is the full presentation of the stated
free-by-free group.  Relations (B5.5) eliminate six of the eleven core
letters, and (B5.4) identifies the remaining presentation with (B5.6).
This proves completeness of the twelve core relators.

Keep the core tree, then greedily add edges in the order of B.22--B.24
whenever they join distinct components.  The stages have
\[
\begin{array}{c|rrrr}
 &\text{core}&\text{first tier}&\text{second tier}&\text{full}\\ \hline
 \text{vertices}&50&182&422&542\\
 \text{edges}&60&360&1020&1620\\
 \text{tree edges}&49&181&421&541\\
 \text{free cycles}&11&179&599&1079\\
 \text{relations}&12&180&600&1080.
\end{array}
\]
At the last stage add the $120$ paths with pin choice zero.  Their new
targets are distinct, so they add every $W$ vertex once and create no
cycle.  This gives exactly the specified tree.

Let $T_v$ be its rooted path and let $e:s\to t$ be a non-tree edge.
The based loop is $T_s eT_t^{-1}$ in traversal order.  Combing its
verified braid gives a word $w_e(a,b,x,y,z)$; substitution (B5.4)
gives its explicit word $W_e(g_1,\ldots,g_{11})$ in core cycles.
For each of the $1068$ non-core non-tree edges impose
\[
                         T_s eT_t^{-1}=W_e.                 \tag{B5.7}
\]
The JSON file stores both sides as full oriented edge words, together
with the braid word and its five-generator normal form.  All $1068$
equalities, and all twelve expanded core equalities, have been checked
in the faithful Artin representation.

Completeness is a Tietze argument with explicit inverse maps.  Collapse
the chosen tree in the free groupoid.  It leaves $1079$ free based
letters.  Relations (B5.7) each eliminate its own one new letter and
use only the eleven core letters on the right.  What remains is exactly
the twelve-relator core presentation just proved.  The inverse sends
the five standard moduli generators to (B5.4), and tree transports
extend it to every object.  This proves the claimed equivalence on
all Hom-sets and its natural inverse, as well as the count
$1068+12=1080$.
\end{proof}

\begin{result}{B.5.1}{Proved}{Exact transport to other path conventions}
Fix the objects, the core paths and the spanning tree of B.5.
Relative path assignments on the other arrows form a torsor for
$\Gamma_{0,5}^{1068}$.  For each such assignment there is an explicit
invertible change of typed generators carrying the $1080$ relations
of B.5 to a complete relation system for that assignment.
The incidence data, fractional-twist arrows and endpoint plumbing
germs alone do not determine the other path classes.
\end{result}

\begin{proof}
Let $F$ be the realization proved above and $Q$ a path assignment
agreeing on the core and tree.  The discrepancy for an unfixed edge
$e:s\to t$ is the based loop
$\delta_e=[T_sQ(e)F(e)^{-1}T_s^{-1}]$ in $\Gamma_{0,5}$.
Choose its unique free-by-free normal form, then (B5.4) gives a
specified core word $\kappa(\delta_e)$.  Put
$L_e=T_s^{-1}\kappa(\delta_e)T_s$, a loop at the source $s$.
The free-groupoid automorphism
\[
 \Phi(e)=L_e e\quad\text{on unfixed edges},\qquad
 \Phi(e)=e\quad\text{on the core and tree}
\]
has inverse obtained by replacing $L_e$ by $L_e^{-1}$, because every
$L_e$ uses only fixed edges.  The equality $Q=F\Phi$ follows directly
from the discrepancy definition.  Consequently the complete relation
list for $Q$ is $\Phi^{-1}(\mathcal R)$, where $\mathcal R$ is the
computed list of B.5.  This proves the exact comparison, including its
inverse and completeness, for all choices rather than equating their
arrow names.

For nonuniqueness, before each of the $120$ composite $f_{\vec Q,a}$
arrows insert the positive meridian of $q=0$ in
$V(q)=(\infty,-q,q,-1,1)$, based at $q=i\epsilon$.
It remains within the same four-point double-cover family.  After
normalization its braid is the pair-meridian, conjugate to
$\sigma_2^2=A_{23}$, whose class in
$H_1(M_{0,5};\mathbb Z)=\mathbb Z^5$ is nonzero.
Relabel this construction by all permutations.  The stabilizer
$(12)(34)$ acts by the constant projective map $z\mapsto-z$ on this
family, so it fixes the meridian path and the modification is
well-defined on source cosets.  The resulting assignment is
$\mathfrak S_5$-equivariant and changes exactly those $120$ paths.
Insert the loop after a short fixed initial segment and before a short
fixed terminal segment to keep both endpoint germs literally unchanged.
All fractional-twist arrows, endpoints and incidences are unchanged,
while the nonzero homology discrepancy proves distinct relative path
classes.  Thus a further marking-forgetting equation or direct path
convention must be compared by its actual discrepancy words; it cannot
be inferred from those retained incidence data.

DELTA\_ON\_GT B.33 correctly identified this requirement.  B.5 now
provides the complete computed fundamental relation system for the
specified realization of the diagram in cote~148, scans~31--33.
B.5.1 gives the exact morphism for transporting it to any further
specified path realization.  The archive's unrecorded manuscript
choices are not silently assigned zero discrepancy.
\end{proof}

\paragraph{The later manuscript marking operations.}
We also opened cote~148, folios~51--56 (PDF scans~52--57),
including the three erasing operations, mixed erasure, and
uniform additional marking \cite[folios~51--56]{GrothendieckCote148}.
They define actual functors on surfaces equipped with their
named cutting circles and marking subsets.  The uniform
operations give the displayed maps of cyclic marking torsors.
To identify their arrows with B.5, the marked surfaces must
still be assigned to the chosen affine configurations, and
the resulting transports compared by the words of B.5.1.
The read pages contain no such computed transport words.
The retained source audit records the exact legible passages
and leaves the heavily corrected marginal text unused.
